\documentclass[11pt,a4paper]{article}
\usepackage[T1]{fontenc}
\usepackage[margin=2.5cm]{geometry}
\usepackage{amsmath}
\pagestyle{empty}
% Pole "Abstract" w EasyChair, GCMM 2026, submission 59.
% SAM ABSTRAKT -- bez tytulu, slow kluczowych, autorow i afiliacji.
% Konferencja nie wymaga tam tytulu ani slow kluczowych, a zgloszenie jest slepe.
% Tekst znak w znak jak abstrakt w main.tex: 249 slow (limit strony 150-250).
\begin{document}
\noindent
The permutation flow shop problem (PFSP) is a canonical NP-hard
scheduling problem whose local-search solvers spend most of their
effort evaluating neighborhood moves. Recent work has offloaded this
evaluation to quantum annealers by casting neighborhoods as QUBO
problems, yet the gate-model paradigm, despite its flexible
connectivity, remains largely unexplored for scheduling. This paper
introduces a gate-model framework that evaluates PFSP neighborhoods with
the Quantum Approximate Optimization Algorithm (QAOA) embedded in
classical Iterated Local Search and Simulated Annealing metaheuristics.
We reuse the swap-selection QUBO encoding from the annealing
line, so a single problem matrix drives either an annealer or a QAOA
circuit, and we adopt a fixed-angle QAOA whose variational angles are
pre-optimized offline; on hardware each move then costs a single device
submission, which keeps the method viable under NISQ constraints. We
formulate all four neighborhood structures (Adjacent, Fibonacci,
Dynasearch, Motzkin) and analyze how their conflict-graph density
governs circuit depth once the QUBO is compiled onto superconducting
hardware, from a heavy-hex lattice to a resonator-coupled star. Depth
assumes the role that embedding capacity plays on annealers and produces
the same feasibility ordering for a different, compilation-grounded
reason, and windowed decomposition tames the dense neighborhoods.
Measurements on three processors separate two regimes: neighborhoods
whose circuits stay short return the same makespan on a 156-qubit
heavy-hex device and on a 16-qubit resonator device, instance by
instance, while the deepest circuits shift between coupling maps. The
paper sets out the evaluation methodology, reports this execution, and
states the open challenges that remain.
\end{document}
